% !TEX root = ../main.tex
% !TEX program = lualatex
\documentclass[../main.tex]{subfiles}
\IfSubfilesClassLoaded{\externaldocument{\subfix{../../build/main}}}{}
% =============================================================================
% CHAPTER 27: CIVPERS (CIVILIAN PERSONNEL)
% DESCRIPTION: CIVPERS overview, workforce shaping, \ac{eeo}, employee and labor relations.
% =============================================================================
\AtBeginDocument{\chapter{EDO Study Guide}}{}
\begin{document}
\IfSubfilesClassLoaded{\chapter{EDO Study Guide}}{}
%====================
% Contents here
%====================

\section{Civilian Personnel Fundamentals}

\subsection{Why CIVPERS management matters}
\begin{description}
  \item[\textbf{Workforce weight.}] \ac{civpers} form the majority of the engineering and logistics talent at \acp{syscom} and \acp{nwc}; how you engage and develop them directly sets Fleet readiness outcomes~\autocite{edo-6-1-1-civpers-overview-2024}.
  \item[\textbf{\ac{hro} lanes.}] Core services are staffing, classification/compensation, employee relations, and labor relations; pull them in early for hiring plans, performance actions, and workplace changes~\autocite{edo-6-1-1-civpers-overview-2024}.
  \item[\textbf{Guidance stack.}] Primary references flow from Title 5, \ac{opm} guidance, Department of the Navy policy, and local agreements; document assumptions and keep counsel/\ac{hro} informed to avoid missteps~\autocite{edo-6-1-1-civpers-overview-2024}.
\end{description}

\subsection{Workforce shaping and classification}
\begin{description}
  \item[\textbf{Pay systems.}] Know the differences between \ac{gs}, \ac{wg}, pay-band systems, and \ac{ses} tiers; align the billet to the right system before recruiting~\autocite{edo-6-1-2-workforce-shaping-2025}.
  \item[\textbf{Classification and \ac{flsa}.}] The \ac{pd} drives pay plan, series, and grade/band; verify \ac{flsa} exempt/non-exempt status and ensure duties match the description to avoid backpay claims~\autocite{edo-6-1-2-workforce-shaping-2025}.
  \item[\textbf{Filling vacancies.}] Compare reassignment/merit promotion, \ac{dha}, or competitive hiring; honor mandatory placement programs and capture recruiting timelines up front~\autocite{edo-6-1-2-workforce-shaping-2025}.
  \item[\textbf{Downsizing levers.}] Sequence \ac{rif}, \ac{sip}, \ac{vera}, and furlough options carefully; pair with workforce development to retain critical skills while meeting budget constraints~\autocite{edo-6-1-2-workforce-shaping-2025}.
\end{description}

\subsection{Equal Employment Opportunity anchors}
\begin{description}
  \item[\textbf{Bases and theories.}] The ten protected bases (race, color, religion, retaliation, sex, national origin, equal pay, age, disability, genetic information) and theories (disparate treatment, disparate impact, harassment, reprisal) frame every \ac{eeo} review~\autocite{edo-6-1-3-eeo-2025}.
  \item[\textbf{Process discipline.}] Preserve timelines: timely contact with a \ac{eeo} counselor, alternative dispute resolution where appropriate, investigation, and decision/hearing; supervisors must avoid reprisal and secure records~\autocite{edo-6-1-3-eeo-2025}.
  \item[\textbf{Harassment prevention.}] Distinguish hostile environment from quid pro quo scenarios, act on any notice (verbal or written), and coordinate with \ac{hro}/\ac{eeoc} liaisons for training and documentation~\autocite{edo-6-1-3-eeo-2025}.
\end{description}

\subsection{Employee relations essentials}
\begin{description}
  \item[\textbf{Work schedules and leave.}] Use standard, flexible, or \ac{cws} structures thoughtfully; apply annual/sick leave and leave without pay rules consistently and reinforce timekeeping discipline to avoid \ac{awol} issues~\autocite{edo-6-1-4-employee-relations-2025}.
  \item[\textbf{Performance cycle.}] Set clear standards early, conduct midpoint feedback, and close out on time; tailor awards (time-off, cash, quality step increases) to reinforce desired behaviors~\autocite{edo-6-1-4-employee-relations-2025}.
  \item[\textbf{Conduct v.s.\ performance.}] Separate misconduct (e.g., \ac{awol}, falsification) from poor performance; use counseling and a documented \ac{pip} before escalating to adverse action, and keep counsel/\ac{hro} informed~\autocite{edo-6-1-4-employee-relations-2025}.
\end{description}

\subsection{Labor relations touchpoints}
\begin{description}
  \item[\textbf{Bargaining unit scope.}] Confirm who is in the bargaining unit (non-supervisory, non-confidential employees) and what the collective bargaining agreement covers; track official time and steward access obligations~\autocite{edo-6-1-5-labor-relations-2025}.
  \item[\textbf{Management rights v.s.\ bargaining.}] Management retains the right to assign work, direct employees, and determine budget; bargain procedures and impact when changing conditions of employment and document notice-to-union steps~\autocite{edo-6-1-5-labor-relations-2025}.
  \item[\textbf{Past practice and change-of-command.}] There is no grace period for a new \ac{oic}; changes to established bargaining-unit conditions must satisfy the \ac{flra} past-practice test (clear, consistent, long-standing, mutually accepted) and trigger notice-and-bargain steps before implementation. Use the required within-90-day \ac{cca} to surface desired changes early with counsel/\ac{hro} to avoid implied acceptance of prior practices~\autocite{FLRA-IRS-3-FLRA-513,OPNAVINST-5354-1H}.
  \item[\textbf{Meetings and disputes.}] Provide the union a chance to attend formal discussions and Weingarten meetings; log official time and resolve \acp{ulp} through grievance procedures or \ac{flra} channels~\autocite{edo-6-1-5-labor-relations-2025}.
\end{description}

%====================
% End of file
%====================
\IfSubfilesClassLoaded{
  \RenewDocumentCommand{\entryname}{}{\textbf{\color{Modern} Acronym}}
  \RenewDocumentCommand{\descriptionname}{}{\textbf{\color{Modern} Definition}}
  \printnoidxglossary[
    type=\acronymtype,
    title=Acronyms,
    style=long-booktabs
  ]}{}
\end{document}
